\notechapter{N-20220811}{Inequality Exercises: Constraints, Equality, and Vector Geometry}

The central skill in an inequality problem is finding the correct representation of the expression. Fractions can be rewritten as a Cauchy--Schwarz inner product; a radical can encode a Euclidean distance; a quadratic constraint defines an ellipsoid or sphere. The standard methods and equality conditions follow Hardy, Littlewood and Pólya's \emph{Inequalities} and Strang's MIT 18.06 linear algebra material.\footnote{\url{https://web.mit.edu/18.06/www/index.shtml}}

\section{Fraction inequalities and their domains}

For positive denominators \(u,v\), Cauchy--Schwarz gives
\[
 \frac{a^2}{u}+\frac{b^2}{v}
 \ge\frac{(a+b)^2}{u+v}.
\]
To see equality, write
\[
 \frac{a^2}{u}+\frac{b^2}{v}-\frac{(a+b)^2}{u+v}
 =\frac{(av-bu)^2}{uv(u+v)}.
\]
Thus equality occurs exactly when \(a/u=b/v\).

The original expression
\[
 F(x,y)=\frac4{4-x^2}+\frac9{9-y^2}
\]
can be bounded by the same mechanism \emph{only} when \(|x|<2\) and \(|y|<3\). Under these conditions,
\[
 F(x,y)\ge\frac{25}{13-x^2-y^2}.
\]
If in addition \(xy=-1\), then \(x^2+y^2\ge2|xy|=2\), so
\[
 \boxed{F(x,y)\ge\frac{25}{11}.}
\]
However this last bound need not be sharp: equality in Cauchy would require \(2/(4-x^2)=3/(9-y^2)\), and equality in \(x^2+y^2\ge2\) would require \(x^2=y^2=1\), conditions which cannot both hold. It is an instructive valid \emph{estimate}, not necessarily the optimum.

\section{A square-root comparison is Cauchy in disguise}

Let \(a,b,c,d,m,n>0\). Define
\[
 P=\sqrt{ab}+\sqrt{cd},\qquad
 Q=\sqrt{ma+nc}\sqrt{b/m+d/n}.
\]
Apply Cauchy to
\[
 u=(\sqrt{ma},\sqrt{nc}),\qquad
 v=(\sqrt{b/m},\sqrt{d/n}).
\]
Their inner product is \(P\), and \(\|u\|\|v\|=Q\). Hence \(P\le Q\). Equality occurs precisely when \(u\) and \(v\) are positive scalar multiples, equivalently
\[
 \frac{\sqrt{ma}}{\sqrt{b/m}}
 =\frac{\sqrt{nc}}{\sqrt{d/n}}
 \quad\Longleftrightarrow\quad
 \boxed{m^2ad=n^2bc}.
\]
The equality condition is a useful check on the expansion of \(Q^2\).

\section{Maximizing two radicals on a fixed interval}

For \(3\le x\le5\), consider
\[
 f(x)=2\sqrt{x-3}+\sqrt{5-x}.
\]
Cauchy--Schwarz gives
\[
 f(x)^2\le(2^2+1^2)((x-3)+(5-x))=10.
\]
Equality holds when
\[
 (\sqrt{x-3},\sqrt{5-x})=\lambda(2,1)
 \quad\text{with }\lambda\ge0.
\]
Thus \((x-3)/(5-x)=4\), yielding \(x=23/5\). Direct substitution confirms
\[
 \boxed{\max_{3\le x\le5}f(x)=\sqrt{10}.}
\]
This proof establishes both a universal upper bound and a point where it is attained; there is no need to differentiate and then check endpoints separately.

\section{Equality in Cauchy recovers unknown vectors}

Suppose \(u=(a,b,c)\) and \(v=(x,y,z)\) satisfy
\[
 \|u\|^2=25,\quad \|v\|^2=36,\quad \langle u,v\rangle=30.
\]
Cauchy gives \(\langle u,v\rangle\le\|u\|\|v\|=5\cdot6=30\). Equality forces \(u=\lambda v\) for some \(\lambda\ge0\). Comparing norms gives
\[
 u=\frac56v.
\]
Consequently \(a+b+c=\frac56(x+y+z)\). The ratio \((a+b+c)/(x+y+z)\) equals \(5/6\) \emph{only if} the denominator is nonzero. When \(x+y+z=0\), both sums vanish, so their ratio is \(0/0\) and is undefined. The proportionality of the vectors remains valid.

\section{Triangle angles and sharp symmetric bounds}

If \(A,B,C>0\) are the angles of a Euclidean triangle, \(A+B+C=\pi\). Engel's inequality gives
\[
 \frac1A+\frac1B+\frac1C
 \ge\frac{(1+1+1)^2}{A+B+C}
 =\frac9\pi.
\]
Cauchy also gives
\[
 A^2+B^2+C^2\ge\frac{(A+B+C)^2}{3}
 =\frac{\pi^2}{3}.
\]
The equality cases agree: both are attained by the equilateral triangle \(A=B=C=\pi/3\).

\section{Two spheres: subtraction gives the radical plane}

Let \(c_1,c_2\in\mathbb R^3\) and consider spheres
\[
 \|x-c_1\|^2=r_1^2,\qquad \|x-c_2\|^2=r_2^2.
\]
Subtracting their equations cancels \(\|x\|^2\) and leaves
\[
 2(c_2-c_1)\cdot x
 =r_1^2-r_2^2+\|c_2\|^2-\|c_1\|^2.
\]
If \(c_1\ne c_2\), this is a plane, called the radical plane. Depending on how that plane meets either sphere, the intersection is a circle, a single point, or empty.

For a separate constrained optimization, maximize \(w\cdot x\) on \(\|x-c\|=r\). Writing \(x=c+ru\), \(\|u\|=1\), Cauchy gives
\[
 w\cdot x=w\cdot c+r(w\cdot u)
 \le w\cdot c+r\|w\|,
\]
with equality at \(u=w/\|w\|\) when \(w\ne0\). For \(w=(-8,6,-24)\), \(\|w\|=26\). Thus the maximum is \(w\cdot c+26r\), with its location determined explicitly by the coefficient vector.

\section{A product constraint and AM--GM}

Suppose \(a,b,c>0\) and \((1+a)(1+b)(1+c)=8\). Since
\[
 (1+a)(1+b)(1+c)
 \ge\left(1+(abc)^{1/3}\right)^3,
\]
we have \(2\ge1+(abc)^{1/3}\) and therefore \(abc\le1\).

To justify the middle inequality, let \(t=(abc)^{1/3}\). AM--GM gives \(a+b+c\ge3t\) and \(ab+bc+ca\ge3t^2\), so expansion of the left side is at least \(1+3t+3t^2+t^3\). Equality requires \(a=b=c=1\), where \(abc=1\) and the constraint holds. The bound is therefore sharp.

\section{Three rays turn a radical into a triangle}

Consider the inequality, valid for arbitrary real \(x,y,z\),
\[
 \sqrt{x^2+xy+y^2}
 +\sqrt{x^2+xz+z^2}
 \ge\sqrt{y^2+yz+z^2}.
\]
Choose three unit vectors \(e_1,e_2,e_3\) in the plane separated by \(120^\circ\), so \(e_i\cdot e_j=-1/2\) for \(i\ne j\). Put
\[
 A=xe_1,\qquad B=ye_2,\qquad C=ze_3.
\]
Then
\[
 \|A-B\|^2=x^2+xy+y^2,\quad
 \|A-C\|^2=x^2+xz+z^2,\quad
 \|B-C\|^2=y^2+yz+z^2.
\]
The proposed inequality is exactly \(\|A-B\|+\|A-C\|\ge\|B-C\|\). It is the triangle inequality for three explicit points. The three-ray construction is the missing geometric detail: simply identifying each radical separately as a vector norm would not prove the relationship among them.

\section{Quadratic constraints and eigenvalues}

If \(H\) is a real symmetric matrix, the extremal values of \(x^THx\) on the unit sphere equal its smallest and largest eigenvalues. Indeed, choose an orthonormal eigenbasis, write \(x=\sum_i c_ie_i\), and calculate
\[
 \frac{x^THx}{x^Tx}
 =\frac{\sum_i\lambda_i c_i^2}{\sum_i c_i^2}.
\]
The quotient is a weighted average of the eigenvalues and lies between their minimum and maximum. Equality occurs on the corresponding eigenspaces.

This links inequality practice to constrained optimization without replacing concrete proofs by a generic appeal to Lagrange multipliers. When a question actually contains a quadratic form and a sphere constraint, the eigenvalue calculation is its natural solution.
